% GENERATED FILE. Edit canonical lessons, then run npm run generate:books.
% canonical-source-hash: fnv1a64:50403c9e67a9e065
% canonical-lessons: KA-C06-dative-ge, KA-C06-dative-stacking, KA-C06-dative-subject

\chapter{Case Endings and Dative Subjects}
\label{ch:ka-dative-case}

This chapter is generated from the canonical micro-lessons used by Language Ladder.
Each section stays independently resumable and preserves the authored prerequisite order.

\section[-ge]{-\kn{ಗೆ} (-ge) — "to, for"}

\label{lesson:KA-C06-dative-ge}



\begin{quote}
\textbf{Warm-up.} {[}PAUSE 2s{]} Your first \textbf{case ending}. English puts "to" in \textbf{front}, as a separate word. Kannada sticks it on the \textbf{back} — the biggest structural fact about the language, in one suffix.
\end{quote}

\subsection*{You'll want to know: Adding it on}
\noindent
\begin{tabularx}{\linewidth}{@{}>{\raggedright\arraybackslash}X>{\raggedright\arraybackslash}X>{\raggedright\arraybackslash}X@{}}
\toprule
\textbf{word} & \textbf{+ -ge} & \textbf{meaning} \\
\midrule
\kn{ಹೆಸರು} \emph{hesaru} — name & \kn{ಹೆಸರ}\textbf{\kn{ಿಗೆ}} \emph{hesar\textbf{ige}} & to/for the name \\
\kn{ಕೆಲಸ} \emph{kelasa} — work & \kn{ಕೆಲಸ}\textbf{\kn{ಕ್ಕೆ}} \emph{kelasa\textbf{kke}} & for work \\
\bottomrule
\end{tabularx}

Two shapes, one suffix: \textbf{-ige} after some stems, \textbf{-kke} after others. Kannada adjusts it to fit the word in front — but you can still see exactly where the noun stops and the ending starts.

And with "I", which reshapes first:

\begin{quote}
\textbf{\kn{ನಾನು}} \emph{nānu} ("I") $\to$ \textbf{\kn{ನನಗೆ}} \emph{nanage} ("to me")
\end{quote}

\begin{cousinweb}
All four descend from one Proto-Dravidian dative, \emph{\textbf{-k(k)u}}. Look at what became of its \emph{k}:

\noindent
\begin{tabularx}{\linewidth}{@{}>{\raggedright\arraybackslash}X>{\raggedright\arraybackslash}X@{}}
\toprule
\textbf{language} & \textbf{dative} \\
\midrule
Tamil & \emph{-ukku} \\
Telugu & \emph{-ku} / \emph{-ki} \\
Malayalam & \emph{-ikku} / \emph{-inu} \\
\textbf{Kannada} & \emph{\textbf{-ge}} / \emph{-ige} \\
\bottomrule
\end{tabularx}

Kannada is the one that \textbf{voiced} it — \emph{k} $\to$ \emph{g} between vowels — and carried that change into the dative, where its sisters kept the hard consonant. Note the other Kannada form, \emph{\textbf{-kke}}: a \textbf{doubled} consonant resists voicing, so the geminate preserved the original \emph{k}. The two Kannada shapes are the same suffix caught on either side of a sound change.
\end{cousinweb}

\subsection*{Guided Practice}
{[}PAUSE 1s{]}

\begin{itemize}
\raggedright
  \item {[}YOU SAY: "\emph{hesaru} … \emph{hesarige}"{]}
  \item {[}YOU SAY: "\emph{nānu} … \emph{nanage}" — "I" … "to me"{]}
  \item {[}YOU SAY: "\emph{kelasakke}" — "for work," from Ch. 5's \kn{ಕೆಲಸ}{]}
  \item {[}YOU SAY: the family — "\emph{-ukku · -ku · -\textbf{ge}} · -ikku"{]}
\end{itemize}

\begin{tcolorbox}[breakable,colback=teal!4,colframe=teal!35!black,title={Wrap-up recall}]
{[}PAUSE 3s{]} What does \textbf{-\kn{ಗೆ}} mean? ("\textbf{To}" or "\textbf{for}.") Where does it go? (\textbf{On the end} of the noun.) What is "to me"? (\textbf{\kn{ನನಗೆ}} \emph{nanage}.) Why does Kannada have \emph{g} where its sisters have \emph{k}? (Kannada \textbf{voiced} the \emph{k} between vowels; doubled \emph{-kke} resisted and kept the old hard form.)
\end{tcolorbox}
\section[-ge]{-\kn{ಗೆ} — one suffix, one job}

\label{lesson:KA-C06-dative-stacking}



\begin{quote}
\textbf{Warm-up.} {[}PAUSE 2s{]} You can now attach \textbf{-\kn{ಗೆ}} and make \textbf{\kn{ನನಗೆ}} “to me.” Look at what that visible seam tells you about how Kannada builds longer words.
\end{quote}

\begin{grammarlens}[title={Grammar Lens: Stacking versus fusion}]
\noindent
\begin{tabularx}{\linewidth}{@{}>{\raggedright\arraybackslash}X>{\raggedright\arraybackslash}X>{\raggedright\arraybackslash}X@{}}
\toprule
\textbf{} & \textbf{Kannada \textbf{-ge}} & \textbf{a Latin ending like \textbf{-īs}} \\
\midrule
pins down the case? & \textbf{yes} — dative, always & \textbf{no} — dative \textbf{or} ablative \\
carries number? & \textbf{no} — a separate suffix does & yes — plural, baked in \\
carries declension class? & \textbf{no} & yes — first or second \\
can you see the seam? & \textbf{yes}: \emph{hesar} + \emph{ige} & no — one fused lump \\
\bottomrule
\end{tabularx}

Kannada \textbf{stacks} suffixes. Each suffix carries one meaning and sits in a fixed order, so the learner can see where the noun ends and the case marker begins. That is what \textbf{agglutinative} means here: longer words are assembled from pieces whose jobs remain visible.

Latin \textbf{fuses} several jobs into one ending. \emph{-īs} carries case, plural number, and declension class at once—and it does not even settle whether the case is dative or ablative. Kannada \textbf{-\kn{ಗೆ}} carries one unambiguous job: “to/for.”
\end{grammarlens}

\subsection*{Guided Practice}
{[}PAUSE 1s{]}

\begin{itemize}
\raggedright
  \item {[}YOU SAY: \emph{hesar} + \emph{ige} — a visible noun-plus-case seam{]}
  \item {[}YOU SAY: Kannada \emph{-ge} — one job, dative{]}
  \item {[}YOU SAY: Latin \emph{-īs} — case, plural, and declension fused{]}
\end{itemize}

\begin{tcolorbox}[breakable,colback=teal!4,colframe=teal!35!black,title={Wrap-up recall}]
{[}PAUSE 3s{]} Why is Kannada called agglutinative? (\textbf{It stacks visible pieces, each with one job.}) How many jobs can Latin \emph{-īs} fuse? (\textbf{Case, number, and declension class—and its case can still be dative or ablative.})
\end{tcolorbox}
\section[nanage kannaḍa gottu]{\kn{ನನಗೆ} \kn{ಕನ್ನಡ} \kn{ಗೊತ್ತು} — "I know Kannada"}

\label{lesson:KA-C06-dative-subject}



\begin{quote}
\textbf{Warm-up.} {[}PAUSE 2s{]} Chapter 5 gave you \textbf{\kn{ನಾನು} \kn{ಕನ್ನಡ} \kn{ಮಾತನಾಡುತ್ತೇನೆ}} — "\textbf{I} speak Kannada." Now say "I \textbf{know} Kannada." Kannada will not let you use \emph{nānu}.
\end{quote}

\subsection*{You'll want to know: The sentence}
\begin{quote}
\textbf{\kn{ನನಗೆ} \kn{ಕನ್ನಡ} \kn{ಗೊತ್ತು}.} \emph{nanage kannaḍa gottu.} literally: "\textbf{to-me} Kannada {[}is{]} \textbf{known}."
\end{quote}

\noindent
\begin{tabularx}{\linewidth}{@{}>{\raggedright\arraybackslash}X>{\raggedright\arraybackslash}X@{}}
\toprule
\textbf{piece} & \textbf{} \\
\midrule
\textbf{\kn{ನನಗೆ}} \emph{nanage} & "to me" — the dative from the last lesson \\
\textbf{\kn{ಕನ್ನಡ}} \emph{kannaḍa} & Kannada (Ch. 5) \\
\textbf{\kn{ಗೊತ್ತು}} \emph{gottu} & "known" \\
\bottomrule
\end{tabularx}

Two things to notice. \textbf{You are not the subject} — \emph{nānu} is gone, moved into the dative, leaving \emph{kannaḍa} as the unmarked word the sentence is about. And \emph{gottu} isn't really a verb doing an action; it's closer to "\textbf{{[}is{]} known}." So nothing in this sentence is \emph{doing} anything at all.

\begin{grammarlens}[title={Grammar Lens: Why: things that happen \emph{to} you}]
Kannada separates what you \textbf{do} from what \textbf{happens to} you. Speaking is an action, so Ch. 5 used \emph{nānu}. But knowing, liking, wanting, being hungry — these \textbf{arrive}. Kannada marks that by putting the person in the \textbf{dative}.

\noindent
\begin{tabularx}{\linewidth}{@{}>{\raggedright\arraybackslash}X>{\raggedright\arraybackslash}X@{}}
\toprule
\textbf{English frames it as} & \textbf{Kannada frames it as} \\
\midrule
I \textbf{know} Kannada & Kannada is known \textbf{to me} \\
I \textbf{like} it & it is pleasing \textbf{to me} \\
I \textbf{am} cold & cold is \textbf{to me} \\
\bottomrule
\end{tabularx}

English keeps one fossil of the same idea: "\textbf{methinks}" — where \emph{me} is a dative, "it seems \textbf{to me}." Kannada made it a rule.
\end{grammarlens}

\begin{grammarlens}[title={Grammar Lens: All four sisters do it}]
\noindent
\begin{tabularx}{\linewidth}{@{}>{\raggedright\arraybackslash}X>{\raggedright\arraybackslash}X>{\raggedright\arraybackslash}X@{}}
\toprule
\textbf{language} & \textbf{"I know {[}the language{]}"} & \textbf{the dative} \\
\midrule
Kannada & \emph{nanage kannaḍa gottu} & \textbf{-ge} \\
Tamil & \emph{enakku tamiḻ teriyum} & \textbf{-ukku} \\
Telugu & \emph{nāku telugu vaccu} & \textbf{-ku} \\
Malayalam & \emph{enikku malayāḷam aṟiyām} & \textbf{-ikku} \\
\bottomrule
\end{tabularx}

One construction, four languages, four suffixes that are visibly the \textbf{same suffix} — Kannada's \emph{g} being the softened \emph{k} of the others. The Dravidian family showing its bones, exactly as \emph{blanc/bianco/branco} did for Romance.
\end{grammarlens}

\subsection*{Guided Practice}
{[}PAUSE 1s{]}

\begin{itemize}
\raggedright
  \item {[}YOU SAY: "\emph{nanage kannaḍa gottu}"{]}
  \item {[}YOU SAY: the contrast — "\emph{nānu kannaḍa mātanāḍuttēne}" (I speak) … "\emph{nanage kannaḍa gottu}" (known to me){]}
  \item {[}YOU SAY: the literal — "to-me Kannada known"{]}
  \item {[}YOU SAY: the four cousins — "\emph{-ge · -ukku · -ku · -ikku}"{]}
\end{itemize}

\begin{tcolorbox}[breakable,colback=teal!4,colframe=teal!35!black,title={Wrap-up recall}]
{[}PAUSE 3s{]} What does \textbf{\kn{ನನಗೆ} \kn{ಕನ್ನಡ} \kn{ಗೊತ್ತು}} literally say? ("\textbf{To-me} Kannada \textbf{known}.") What has English got that Kannada hasn't? (A nominative "\textbf{I}" — \emph{nānu} is replaced by the dative \emph{nanage}.) Is \emph{gottu} an action verb? (\textbf{No} — it's "{[}is{]} known"; nothing is doing anything.) Why the dative? (Knowing \textbf{happens to} you.) Which English word keeps the same idea? (\textbf{Methinks}.)
\end{tcolorbox}
